\noindent\textbf{(a)} La factorización es \(P(r)=(r-a)[r^2+(2-a)r+1]\). Si \(0<a<4\), con \(\nu=\sqrt{a(4-a)}/2\),
\[y=C_1e^{ax}+e^{(a-2)x/2}(C_2\cos\nu x+C_3\sin\nu x).\]
Los casos especiales son
\[\begin{array}{ll}a=0:&y=C_1+e^{-x}(C_2+C_3x),\\a=4:&y=C_1e^{4x}+e^x(C_2+C_3x),\\a=-\tfrac12:&y=C_1e^{-2x}+e^{-x/2}(C_2+C_3x).\end{array}\]
Para \(a<0\), \(a\ne-1/2\), o \(a>4\),
\[y=C_1e^{ax}+C_2e^{r_+x}+C_3e^{r_-x},\quad r_\pm=\frac{a-2\pm\sqrt{a(a-4)}}2.\]

\noindent\textbf{(b)} Si \(a\ne0,\pm1\),
\[y=C_1e^{ax}+C_2e^{-ax}+C_3e^{x/a}.\]
Si \(a=\pm1\),
\[y=(C_1+C_2x)e^{ax}+C_3e^{-ax}.\]
